
\documentclass{article}
\usepackage{colortbl}
\usepackage{makecell}
\usepackage{multirow}
\usepackage{supertabular}

\begin{document}

\newcounter{utterance}

\centering \large Interaction Transcript for game `ifeval', experiment `ifeval', episode 240 with llama{-}3.1{-}8b{-}instruct.
\vspace{24pt}

{ \footnotesize  \setcounter{utterance}{1}
\setlength{\tabcolsep}{0pt}
\begin{supertabular}{c@{$\;$}|p{.15\linewidth}@{}p{.15\linewidth}p{.15\linewidth}p{.15\linewidth}p{.15\linewidth}p{.15\linewidth}}
    \# & \multicolumn{2}{c}{Player} && \multicolumn{2}{c}{Game Master} \\
    \hline

    \theutterance \stepcounter{utterance}  
    & & & \multicolumn{4}{p{0.6\linewidth}}{
        \cellcolor[rgb]{0.9,0.9,0.9}{
            \makecell[{{p{\linewidth}}}]{
                \texttt{\tiny{[P1$\langle$GM]}}
                \texttt{Write a casual blog post about how the outer solar system is different from the inner solar system, and what that means for the possibility of life. Wrap your entire response with double quotation marks. Your response should contain 17 or more sentences.} \\
            }
        }
    }
     \\ \\

    \theutterance \stepcounter{utterance}  
    & \multicolumn{4}{p{0.6\linewidth}}{
        \cellcolor[rgb]{0.561,0.737,0.561}{
            \makecell[{{p{\linewidth}}}]{
                \texttt{\tiny{[P1$\rangle$GM]}}
                \texttt{"Exploring the vastness of our solar system, we often find ourselves comparing the inner and outer regions. While the inner solar system is home to the rocky planets like Earth, Mercury, and Mars, the outer solar system is a realm of gas giants and icy worlds. The most notable difference between the two is the temperature and composition of the planets.} \\
\\ 
\texttt{In the inner solar system, the temperatures are relatively warm, with the average temperature on Earth being around 15°C (59°F). This warmth is due to the proximity to the sun, which provides the necessary energy for liquid water to exist. Liquid water is essential for life as we know it, making the inner solar system a prime candidate for hosting life. The rocky planets in the inner solar system also have solid surfaces, which provide a stable environment for life to thrive.} \\
\\ 
\texttt{On the other hand, the outer solar system is a much colder and more distant region. The gas giants like Jupiter and Saturn are massive balls of gas, with temperatures ranging from {-}150°C to {-}200°C ({-}238°F to {-}330°F). The icy worlds like Uranus and Neptune are even colder, with temperatures as low as {-}200°C ({-}330°F). The lack of warmth and the presence of intense radiation from the sun make the outer solar system a challenging environment for life to exist.} \\
\\ 
\texttt{However, the outer solar system is not entirely devoid of life. The moons of the gas giants, such as Europa and Enceladus, have subsurface oceans that are thought to be in contact with rock. This interaction between water and rock could provide the necessary energy and nutrients for life to thrive. The presence of organic molecules, which are the building blocks of life, has also been detected in the outer solar system.} \\
\\ 
\texttt{The discovery of exoplanets in the outer solar system has also sparked interest in the possibility of life existing beyond our own solar system. The detection of exoplanets with conditions similar to those of Earth has led scientists to speculate about the possibility of life existing elsewhere in the universe. While the outer solar system is a challenging environment for life, it is not impossible.} \\
\\ 
\texttt{In conclusion, the outer solar system is a vastly different region from the inner solar system, with colder temperatures and a more distant environment. While the inner solar system is a prime candidate for hosting life, the outer solar system is not entirely devoid of life. The discovery of subsurface oceans and organic molecules in the outer solar system has sparked interest in the possibility of life existing beyond our own solar system. As we continue to explore the outer solar system, we may uncover new and exciting discoveries that challenge our understanding of life in the universe."} \\
            }
        }
    }
    & & \\ \\

    \theutterance \stepcounter{utterance}  
    & & & \multicolumn{2}{p{0.3\linewidth}}{
        \cellcolor[rgb]{0.9,0.9,0.9}{
            \makecell[{{p{\linewidth}}}]{
                \texttt{\tiny{[GM$|$GM]}}
                {'startend:quotation': {}, 'length_constraints:number_sentences': {'relation': 'at least', 'num_sentences': 17}}
            }
        }
    }
    & & \\ \\

    \theutterance \stepcounter{utterance}  
    & & & \multicolumn{2}{p{0.3\linewidth}}{
        \cellcolor[rgb]{0.9,0.9,0.9}{
            \makecell[{{p{\linewidth}}}]{
                \texttt{\tiny{[GM$|$GM]}}
                \texttt{game\_result = WIN} \\
            }
        }
    }
    & & \\ \\

\end{supertabular}
}

\end{document}
